% 5.2 230 个空间群的单值表示
\section{230 个空间群的单值表示}

在编制所有空间群的表示表时，我们用的是 3.7 节详述、3.8 节以立方密堆积结构（$Fm3m$, $O^{5}_{h}$）与金刚石结构（$Fd3m$, $O^{7}_{h}$）为例演示过的理论。我们要确定的是 $\mathbf{G}^{\mathbf{k}}$——波矢 $\mathbf{k}$ 的小群——对相应布里渊区基本域中所有波矢的表示。当然 $\mathbf{G}^{\mathbf{k}}$ 是无限群。对对称点，问题简化为求 Herring 小群 ${}^{H}\mathbf{G}^{\mathbf{k}} = \mathbf{G}^{\mathbf{k}}/\mathbf{T}^{\mathbf{k}}$ 的表示，其中 $\mathbf{T}^{\mathbf{k}}$ 是布拉维格子那些满足 $\exp(-\mathrm{i}\mathbf{k} \cdot \mathbf{t}_{i}) = +1$ 的对称操作 $\{E \mid \mathbf{t}_{i}\}$ 构成的群；$\mathbf{T}^{\mathbf{k}}$ 是布拉维格子全部平移对称操作群 $\mathbf{T}$ 的子群。群 ${}^{H}\mathbf{G}^{\mathbf{k}}$ 的元素可以由表 3.6 与 3.7 确定，然后通过直接相乘 Seitz 空间群符号——每当 $\exp(-\mathrm{i}\mathbf{k} \cdot \mathbf{t}_{i}) = +1$ 时把 $\{E \mid \mathbf{t}_{i}\}$ 换成 $\{E \mid \mathbf{0}\}$——得到 ${}^{H}\mathbf{G}^{\mathbf{k}}$ 的群乘法表。我们称 ${}^{H}\mathbf{G}^{\mathbf{k}}$ 为 Herring 小群，因为这种处理对称点的方法最初由 Herring (1942) 用于确定六角密堆积与金刚石结构的空间群特征标表。这个群随后可以辨认为表 5.1 中的抽象群 $\mathbf{G}^{n}_{|\mathbf{G}|}$ 之一，从而容易辨认可作为小表示（即满足式 (3.4.3) 的表示）的 $\mathbf{G}^{n}_{|\mathbf{G}|}$ 的相应表示。对对称线上的 $\mathbf{k}$，我们必须通过求中央延拓群 $\overline{\mathbf{G}}^{\mathbf{k}*}$ 的射影表示来得到 $\mathbf{G}^{\mathbf{k}}$ 的表示。对此类 $\mathbf{k}$，$\mathbf{T}/\mathbf{T}^{\mathbf{k}}$ 不含有限个元素，因此试图从有限群的表示确定 $\mathbf{G}^{\mathbf{k}}$ 的小表示时，不能有效地使用 $\mathbf{G}^{\mathbf{k}}/\mathbf{T}^{\mathbf{k}}$（它现在是无限群）；为了找到一个合适的有限群，必须改用中央延拓群 $\overline{\mathbf{G}}^{\mathbf{k}*}$。$\overline{\mathbf{G}}^{\mathbf{k}*}$ 元素的群乘法表用乘法规则确定。于是可以把 $\overline{\mathbf{G}}^{\mathbf{k}*}$ 辨认为表 5.1 的某个抽象群，并辨认同它的小表示相应的那些表示。

可以把每个晶体学点群辨认为某个抽象群的一种体现（表 5.2）。表 5.2 的第 3 列给出与该点群的类相对应的抽象群生成元。例如对 $432\ (O)$，星 $A$ 的小陪群为 $32'\ (D'_{3})$，等等。

\begin{figure}[H]
\centering
\includegraphics[width=0.86\textwidth]{c5_52_p299.jpg}
\caption*{\textbf{表 5.2}\ 用抽象群辨认点群（原书第 286 页扫描）。}
\end{figure}
\begin{figure}[H]
\centering
\includegraphics[width=0.86\textwidth]{c5_52_p300.jpg}
\caption*{表 5.2（续）与表 5.3（原书第 287 页扫描）。}
\end{figure}
\begin{figure}[H]
\centering
\includegraphics[width=0.86\textwidth]{c5_52_p301.jpg}
\caption*{表 5.3（续）、表 5.4 与表 5.5（原书第 288 页扫描）。}
\end{figure}
\begin{figure}[H]
\centering
\includegraphics[width=0.86\textwidth]{c5_52_p302.jpg}
\caption*{表 5.5（续）与表 5.6（原书第 289 页扫描）。}
\end{figure}

\begin{figure}[H]
\centering
\includegraphics[width=0.86\textwidth]{c5_52_p299.jpg}
\caption*{原书第 286 页}
\end{figure}
\begin{figure}[H]
\centering
\includegraphics[width=0.86\textwidth]{c5_52_p300.jpg}
\caption*{原书第 287 页}
\end{figure}
\begin{figure}[H]
\centering
\includegraphics[width=0.86\textwidth]{c5_52_p301.jpg}
\caption*{原书第 288 页}
\end{figure}
\begin{figure}[H]
\centering
\includegraphics[width=0.86\textwidth]{c5_52_p302.jpg}
\caption*{原书第 289 页}
\end{figure}
\begin{figure}[H]
\centering
\includegraphics[width=0.86\textwidth]{c5_52_p303.jpg}
\caption*{原书第 290 页}
\end{figure}
\begin{figure}[H]
\centering
\includegraphics[width=0.86\textwidth]{c5_52_p304.jpg}
\caption*{原书第 291 页}
\end{figure}
\begin{figure}[H]
\centering
\includegraphics[width=0.86\textwidth]{c5_52_p305.jpg}
\caption*{原书第 292 页}
\end{figure}
\begin{figure}[H]
\centering
\includegraphics[width=0.86\textwidth]{c5_52_p306.jpg}
\caption*{原书第 293 页}
\end{figure}
\begin{figure}[H]
\centering
\includegraphics[width=0.86\textwidth]{c5_52_p307.jpg}
\caption*{原书第 294 页}
\end{figure}
\begin{figure}[H]
\centering
\includegraphics[width=0.86\textwidth]{c5_52_p308.jpg}
\caption*{原书第 295 页}
\end{figure}
\begin{figure}[H]
\centering
\includegraphics[width=0.86\textwidth]{c5_52_p309.jpg}
\caption*{原书第 296 页}
\end{figure}
\begin{figure}[H]
\centering
\includegraphics[width=0.86\textwidth]{c5_52_p310.jpg}
\caption*{原书第 297 页}
\end{figure}
\begin{figure}[H]
\centering
\includegraphics[width=0.86\textwidth]{c5_52_p311.jpg}
\caption*{原书第 298 页}
\end{figure}
\begin{figure}[H]
\centering
\includegraphics[width=0.86\textwidth]{c5_52_p312.jpg}
\caption*{原书第 299 页}
\end{figure}
\begin{figure}[H]
\centering
\includegraphics[width=0.86\textwidth]{c5_52_p313.jpg}
\caption*{原书第 300 页}
\end{figure}
\begin{figure}[H]
\centering
\includegraphics[width=0.86\textwidth]{c5_52_p314.jpg}
\caption*{原书第 301 页}
\end{figure}
\begin{figure}[H]
\centering
\includegraphics[width=0.86\textwidth]{c5_52_p315.jpg}
\caption*{原书第 302 页}
\end{figure}
\begin{figure}[H]
\centering
\includegraphics[width=0.86\textwidth]{c5_52_p316.jpg}
\caption*{原书第 303 页}
\end{figure}
\begin{figure}[H]
\centering
\includegraphics[width=0.86\textwidth]{c5_52_p317.jpg}
\caption*{原书第 304 页}
\end{figure}
\begin{figure}[H]
\centering
\includegraphics[width=0.86\textwidth]{c5_52_p318.jpg}
\caption*{原书第 305 页}
\end{figure}
\begin{figure}[H]
\centering
\includegraphics[width=0.86\textwidth]{c5_52_p319.jpg}
\caption*{原书第 306 页}
\end{figure}
\begin{figure}[H]
\centering
\includegraphics[width=0.86\textwidth]{c5_52_p320.jpg}
\caption*{原书第 307 页}
\end{figure}
\begin{figure}[H]
\centering
\includegraphics[width=0.86\textwidth]{c5_52_p321.jpg}
\caption*{原书第 308 页}
\end{figure}
\begin{figure}[H]
\centering
\includegraphics[width=0.86\textwidth]{c5_52_p322.jpg}
\caption*{原书第 309 页}
\end{figure}
\begin{figure}[H]
\centering
\includegraphics[width=0.86\textwidth]{c5_52_p323.jpg}
\caption*{原书第 310 页}
\end{figure}
\begin{figure}[H]
\centering
\includegraphics[width=0.86\textwidth]{c5_52_p324.jpg}
\caption*{原书第 311 页}
\end{figure}
\begin{figure}[H]
\centering
\includegraphics[width=0.86\textwidth]{c5_52_p325.jpg}
\caption*{原书第 312 页}
\end{figure}
\begin{figure}[H]
\centering
\includegraphics[width=0.86\textwidth]{c5_52_p326.jpg}
\caption*{原书第 313 页}
\end{figure}
\begin{figure}[H]
\centering
\includegraphics[width=0.86\textwidth]{c5_52_p327.jpg}
\caption*{原书第 314 页}
\end{figure}
\begin{figure}[H]
\centering
\includegraphics[width=0.86\textwidth]{c5_52_p328.jpg}
\caption*{原书第 315 页}
\end{figure}
\begin{figure}[H]
\centering
\includegraphics[width=0.86\textwidth]{c5_52_p329.jpg}
\caption*{原书第 316 页}
\end{figure}
\begin{figure}[H]
\centering
\includegraphics[width=0.86\textwidth]{c5_52_p330.jpg}
\caption*{原书第 317 页}
\end{figure}
\begin{figure}[H]
\centering
\includegraphics[width=0.86\textwidth]{c5_52_p331.jpg}
\caption*{原书第 318 页}
\end{figure}
\begin{figure}[H]
\centering
\includegraphics[width=0.86\textwidth]{c5_52_p332.jpg}
\caption*{原书第 319 页}
\end{figure}
\begin{figure}[H]
\centering
\includegraphics[width=0.86\textwidth]{c5_52_p333.jpg}
\caption*{原书第 320 页}
\end{figure}
\begin{figure}[H]
\centering
\includegraphics[width=0.86\textwidth]{c5_52_p334.jpg}
\caption*{原书第 321 页}
\end{figure}
\begin{figure}[H]
\centering
\includegraphics[width=0.86\textwidth]{c5_52_p335.jpg}
\caption*{原书第 322 页}
\end{figure}
\begin{figure}[H]
\centering
\includegraphics[width=0.86\textwidth]{c5_52_p336.jpg}
\caption*{原书第 323 页}
\end{figure}
\begin{figure}[H]
\centering
\includegraphics[width=0.86\textwidth]{c5_52_p337.jpg}
\caption*{原书第 324 页}
\end{figure}
\begin{figure}[H]
\centering
\includegraphics[width=0.86\textwidth]{c5_52_p338.jpg}
\caption*{原书第 325 页}
\end{figure}
\begin{figure}[H]
\centering
\includegraphics[width=0.86\textwidth]{c5_52_p339.jpg}
\caption*{原书第 326 页}
\end{figure}
\begin{figure}[H]
\centering
\includegraphics[width=0.86\textwidth]{c5_52_p340.jpg}
\caption*{原书第 327 页}
\end{figure}
\begin{figure}[H]
\centering
\includegraphics[width=0.86\textwidth]{c5_52_p341.jpg}
\caption*{原书第 328 页}
\end{figure}
\begin{figure}[H]
\centering
\includegraphics[width=0.86\textwidth]{c5_52_p342.jpg}
\caption*{原书第 329 页}
\end{figure}
\begin{figure}[H]
\centering
\includegraphics[width=0.86\textwidth]{c5_52_p343.jpg}
\caption*{原书第 330 页}
\end{figure}
\begin{figure}[H]
\centering
\includegraphics[width=0.86\textwidth]{c5_52_p344.jpg}
\caption*{原书第 331 页}
\end{figure}
\begin{figure}[H]
\centering
\includegraphics[width=0.86\textwidth]{c5_52_p345.jpg}
\caption*{原书第 332 页}
\end{figure}
\begin{figure}[H]
\centering
\includegraphics[width=0.86\textwidth]{c5_52_p346.jpg}
\caption*{原书第 333 页}
\end{figure}
\begin{figure}[H]
\centering
\includegraphics[width=0.86\textwidth]{c5_52_p347.jpg}
\caption*{原书第 334 页}
\end{figure}
\begin{figure}[H]
\centering
\includegraphics[width=0.86\textwidth]{c5_52_p348.jpg}
\caption*{原书第 335 页}
\end{figure}
\begin{figure}[H]
\centering
\includegraphics[width=0.86\textwidth]{c5_52_p349.jpg}
\caption*{原书第 336 页}
\end{figure}
\begin{figure}[H]
\centering
\includegraphics[width=0.86\textwidth]{c5_52_p350.jpg}
\caption*{原书第 337 页}
\end{figure}
\begin{figure}[H]
\centering
\includegraphics[width=0.86\textwidth]{c5_52_p351.jpg}
\caption*{原书第 338 页}
\end{figure}
\begin{figure}[H]
\centering
\includegraphics[width=0.86\textwidth]{c5_52_p352.jpg}
\caption*{原书第 339 页}
\end{figure}
\begin{figure}[H]
\centering
\includegraphics[width=0.86\textwidth]{c5_52_p353.jpg}
\caption*{原书第 340 页}
\end{figure}
\begin{figure}[H]
\centering
\includegraphics[width=0.86\textwidth]{c5_52_p354.jpg}
\caption*{原书第 341 页}
\end{figure}
\begin{figure}[H]
\centering
\includegraphics[width=0.86\textwidth]{c5_52_p355.jpg}
\caption*{原书第 342 页}
\end{figure}
\begin{figure}[H]
\centering
\includegraphics[width=0.86\textwidth]{c5_52_p356.jpg}
\caption*{原书第 343 页}
\end{figure}
\begin{figure}[H]
\centering
\includegraphics[width=0.86\textwidth]{c5_52_p357.jpg}
\caption*{原书第 344 页}
\end{figure}
\begin{figure}[H]
\centering
\includegraphics[width=0.86\textwidth]{c5_52_p358.jpg}
\caption*{原书第 345 页}
\end{figure}
\begin{figure}[H]
\centering
\includegraphics[width=0.86\textwidth]{c5_52_p359.jpg}
\caption*{原书第 346 页}
\end{figure}
\begin{figure}[H]
\centering
\includegraphics[width=0.86\textwidth]{c5_52_p360.jpg}
\caption*{原书第 347 页}
\end{figure}
\begin{figure}[H]
\centering
\includegraphics[width=0.86\textwidth]{c5_52_p361.jpg}
\caption*{原书第 348 页}
\end{figure}
\begin{figure}[H]
\centering
\includegraphics[width=0.86\textwidth]{c5_52_p362.jpg}
\caption*{原书第 349 页}
\end{figure}
\begin{figure}[H]
\centering
\includegraphics[width=0.86\textwidth]{c5_52_p363.jpg}
\caption*{原书第 350 页}
\end{figure}
\begin{figure}[H]
\centering
\includegraphics[width=0.86\textwidth]{c5_52_p364.jpg}
\caption*{原书第 351 页}
\end{figure}
\begin{figure}[H]
\centering
\includegraphics[width=0.86\textwidth]{c5_52_p365.jpg}
\caption*{原书第 352 页}
\end{figure}
\begin{figure}[H]
\centering
\includegraphics[width=0.86\textwidth]{c5_52_p366.jpg}
\caption*{原书第 353 页}
\end{figure}
\begin{figure}[H]
\centering
\includegraphics[width=0.86\textwidth]{c5_52_p367.jpg}
\caption*{原书第 354 页}
\end{figure}
\begin{figure}[H]
\centering
\includegraphics[width=0.86\textwidth]{c5_52_p368.jpg}
\caption*{原书第 355 页}
\end{figure}
\begin{figure}[H]
\centering
\includegraphics[width=0.86\textwidth]{c5_52_p369.jpg}
\caption*{原书第 356 页}
\end{figure}
\begin{figure}[H]
\centering
\includegraphics[width=0.86\textwidth]{c5_52_p370.jpg}
\caption*{原书第 357 页}
\end{figure}
\begin{figure}[H]
\centering
\includegraphics[width=0.86\textwidth]{c5_52_p371.jpg}
\caption*{原书第 358 页}
\end{figure}
\begin{figure}[H]
\centering
\includegraphics[width=0.86\textwidth]{c5_52_p372.jpg}
\caption*{原书第 359 页}
\end{figure}
\begin{figure}[H]
\centering
\includegraphics[width=0.86\textwidth]{c5_52_p373.jpg}
\caption*{原书第 360 页}
\end{figure}
\begin{figure}[H]
\centering
\includegraphics[width=0.86\textwidth]{c5_52_p374.jpg}
\caption*{原书第 361 页}
\end{figure}
\begin{figure}[H]
\centering
\includegraphics[width=0.86\textwidth]{c5_52_p375.jpg}
\caption*{原书第 362 页}
\end{figure}
\begin{figure}[H]
\centering
\includegraphics[width=0.86\textwidth]{c5_52_p376.jpg}
\caption*{原书第 363 页}
\end{figure}
\begin{figure}[H]
\centering
\includegraphics[width=0.86\textwidth]{c5_52_p377.jpg}
\caption*{原书第 364 页}
\end{figure}
\begin{figure}[H]
\centering
\includegraphics[width=0.86\textwidth]{c5_52_p378.jpg}
\caption*{原书第 365 页}
\end{figure}
\begin{figure}[H]
\centering
\includegraphics[width=0.86\textwidth]{c5_52_p379.jpg}
\caption*{原书第 366 页}
\end{figure}
\begin{figure}[H]
\centering
\includegraphics[width=0.86\textwidth]{c5_52_p380.jpg}
\caption*{原书第 367 页}
\end{figure}
\begin{figure}[H]
\centering
\includegraphics[width=0.86\textwidth]{c5_52_p381.jpg}
\caption*{原书第 368 页}
\end{figure}
\begin{figure}[H]
\centering
\includegraphics[width=0.86\textwidth]{c5_52_p382.jpg}
\caption*{原书第 369 页}
\end{figure}
\begin{figure}[H]
\centering
\includegraphics[width=0.86\textwidth]{c5_52_p383.jpg}
\caption*{原书第 370 页}
\end{figure}
\begin{figure}[H]
\centering
\includegraphics[width=0.86\textwidth]{c5_52_p384.jpg}
\caption*{原书第 371 页}
\end{figure}
\begin{figure}[H]
\centering
\includegraphics[width=0.86\textwidth]{c5_52_p385.jpg}
\caption*{原书第 372 页}
\end{figure}
\begin{figure}[H]
\centering
\includegraphics[width=0.86\textwidth]{c5_52_p386.jpg}
\caption*{原书第 373 页}
\end{figure}
\begin{figure}[H]
\centering
\includegraphics[width=0.86\textwidth]{c5_52_p387.jpg}
\caption*{原书第 374 页}
\end{figure}
\begin{figure}[H]
\centering
\includegraphics[width=0.86\textwidth]{c5_52_p388.jpg}
\caption*{原书第 375 页}
\end{figure}
\begin{figure}[H]
\centering
\includegraphics[width=0.86\textwidth]{c5_52_p389.jpg}
\caption*{原书第 376 页}
\end{figure}
\begin{figure}[H]
\centering
\includegraphics[width=0.86\textwidth]{c5_52_p390.jpg}
\caption*{原书第 377 页}
\end{figure}
\begin{figure}[H]
\centering
\includegraphics[width=0.86\textwidth]{c5_52_p391.jpg}
\caption*{原书第 378 页}
\end{figure}
\begin{figure}[H]
\centering
\includegraphics[width=0.86\textwidth]{c5_52_p392.jpg}
\caption*{原书第 379 页}
\end{figure}
\begin{figure}[H]
\centering
\includegraphics[width=0.86\textwidth]{c5_52_p393.jpg}
\caption*{原书第 380 页}
\end{figure}
\begin{figure}[H]
\centering
\includegraphics[width=0.86\textwidth]{c5_52_p394.jpg}
\caption*{原书第 381 页}
\end{figure}
\begin{figure}[H]
\centering
\includegraphics[width=0.86\textwidth]{c5_52_p395.jpg}
\caption*{原书第 382 页}
\end{figure}
\begin{figure}[H]
\centering
\includegraphics[width=0.86\textwidth]{c5_52_p396.jpg}
\caption*{原书第 383 页}
\end{figure}
\begin{figure}[H]
\centering
\includegraphics[width=0.86\textwidth]{c5_52_p397.jpg}
\caption*{原书第 384 页}
\end{figure}
\begin{figure}[H]
\centering
\includegraphics[width=0.86\textwidth]{c5_52_p398.jpg}
\caption*{原书第 385 页}
\end{figure}
\begin{figure}[H]
\centering
\includegraphics[width=0.86\textwidth]{c5_52_p399.jpg}
\caption*{原书第 386 页}
\end{figure}
\begin{figure}[H]
\centering
\includegraphics[width=0.86\textwidth]{c5_52_p400.jpg}
\caption*{原书第 387 页}
\end{figure}
\begin{figure}[H]
\centering
\includegraphics[width=0.86\textwidth]{c5_52_p401.jpg}
\caption*{原书第 388 页}
\end{figure}
